

\documentclass[12pt,pdflatex,authoryear]{elsarticle} 
\usepackage[pdftex,hyperref,x11names]{xcolor}
% load the hyper-references package and set document info
\usepackage[pdftex]{hyperref}

% Generate some fake text
\usepackage{blindtext}

% layout control
\usepackage{geometry}
\geometry{verbose,tmargin=1.25in,bmargin=1.25in,lmargin=1.1in,rmargin=1.1in}
\usepackage{parallel}
\usepackage{parcolumns}
\usepackage{fancyhdr}

% math typesetting
\usepackage{array}
\usepackage{amsmath}
\usepackage{amssymb}
\usepackage{amsfonts}
\usepackage{relsize}
\usepackage{mathtools}
\usepackage{bm}
\usepackage[%
decimalsymbol=.,
digitsep=fullstop
]{siunitx}

\usepackage[section]{placeins}

% tables
\usepackage{tabularx}
\usepackage{booktabs}
\usepackage{multicol}
\usepackage{multirow}
\usepackage{longtable}

% to adapt caption style
\usepackage[font={small},labelfont=bf]{caption}

% footnotes at bottom
\usepackage[bottom]{footmisc}
   \renewcommand{\footnotelayout}{\doublespacing} % set spacing in footnotes
   \newlength{\myfootnotesep}
   \setlength{\myfootnotesep}{\baselineskip}
   \addtolength{\myfootnotesep}{-\footnotesep}
   \setlength{\footnotesep}{\myfootnotesep} % set spacing between footnotes

% to change enumeration symbols begin{enumerate}[(a)]
\usepackage{enumerate}

% to make enumerations and itemizations within paragraphs or
% lines. f.i. begin{inparaenum} for (a) is (b) and (c)
\usepackage{paralist}

% graphics stuff
\usepackage{subfig}
\usepackage{graphicx}
\usepackage[space]{grffile} % allows us to specify directories that have spaces
\usepackage{placeins} % prevents floats from moving past a \FloatBarrier
\usepackage{tikz}

% sideway figures
\usepackage{rotating}
\usepackage{lscape}

% Spacing
\usepackage[doublespacing]{setspace}

% remove silly elsevier preprint note
\makeatletter
\def\ps@pprintTitle{%
 \let\@oddhead\@empty
 \let\@evenhead\@empty
 \def\@oddfoot{}%
 \let\@evenfoot\@oddfoot}

\def\input@path{
	{/Users/janus829/Dropbox/Research/conflictEvolution/graphics/}, 	{/Users/s7m/Dropbox/Research/conflictEvolution/graphics/}, 	{/Users/cassydorff/Dropbox/Research/nothingbutnet/conflictEvolution/graphics/}
	       {/Users/maxgallop/Dropbox/conflictEvolutionMex/graphics/} 
}

\graphicspath{
	{/Users/janus829/Dropbox/Research/conflictEvolution/graphics/}, 	{/Users/s7m/Dropbox/Research/conflictEvolution/graphics/}, 	{/Users/cassydorff/Dropbox/Research/nothingbutnet/conflictEvolution/graphics/}
	        {/Users/maxgallop/Dropbox/conflictEvolutionMex/graphics/}  
}

\makeatother

\begin{document}



\renewcommand{\thefigure}{A\arabic{figure}}
\setcounter{figure}{0}
\renewcommand{\thetable}{A.\arabic{table}}
\setcounter{table}{0}
\renewcommand{\thesection}{A.\arabic{section}}
\setcounter{section}{0}

\large{}
\section{\textbf{Appendix: Network Dynamics in Nigeria}}

\subsection{Trace plots for parameter estimates}

Figure~\ref{betaTrace} below shows trace plots for the parameters summarized in Figure 4 of the manuscript.

\begin{figure}[ht]
	\centering
	\includegraphics[width=.88\textwidth]{betaTrace}
	\caption{Trace plots for exogenous parameters.}
	\label{betaTrace}
	\end{figure}
\FloatBarrier
\clearpage

\subsection{Comparison with GLM Estimates}

Figure~\ref{ameglm} compares parameter estimates returned from the AME and GLM frameworks.

\begin{figure}[ht]
	\centering
	\includegraphics[width=.7\textwidth]{ame_v_glm}
	\caption{Comparison of AME and GLM estimates for base model.}
	\label{ameglm}
	\end{figure}
\FloatBarrier
\clearpage

\subsection{Trace plots for parameter estimates when using fatality threshold}

 Figure~\ref{betaTrace_fatalities} below shows trace plots for an alternative formulation of the model we presented in the paper. Specifically, here in defining the dependent variable from ACLED, we only set $y_{ij}=1$ when the corresponding battle between $i$ and $j$ led to at least one fatality. 

\begin{figure}[ht]
	\centering
	\includegraphics[width=.88\textwidth]{betaTrace_fatalities}
	\caption{Trace plots for exogenous parameters.}
	\label{betaTrace_fatalities}
	\end{figure}
\FloatBarrier
\clearpage

\subsection{Trace plots for parameter estimates when setting battles to be symmetric}

 Figure~\ref{betaTrace_symmetric} below shows trace plots for an alternative formulation of the model we presented in the paper. Specifically, here in defining the dependent variable from ACLED, when $y_{ij}=1$ we set $y_{ji}=1$ as well. 

\begin{figure}[ht]
	\centering
	\includegraphics[width=.88\textwidth]{betaTrace_symmetric}
	\caption{Trace plots for exogenous parameters.}
	\label{betaTrace_symmetric}
	\end{figure}
\FloatBarrier
\clearpage

\subsection{Trace plots for parameter estimates when excluding Government actors}

 Figure~\ref{betaTrace_rebelGroup} below shows trace plots for an alternative formulation of the model we presented in the paper. Specifically, here we exclude government actors from the analysis. 

\begin{figure}[ht]
	\centering
	\includegraphics[width=.88\textwidth]{betaTrace_rebelGroupOnly}
	\caption{Trace plots for exogenous parameters.}
	\label{betaTrace_rebelGroup}
	\end{figure}
\FloatBarrier
\clearpage

% \subsection{Additive Effects}

% We depict the sender ($\hat{a}_{i}$) and receiver ($\hat{b}_{j}$) random effects in figure \ref{fig:abEst}. This figure shows which actors are more (and less) violent than would be predicted by just accounting for the exogenous covariates we reviewed in the previous section. We can see, for example, that Boko Haram actually has a random sender effect that is approximately zero, implying that once we account for their high tendency to target civilians and the secular increase in violence that followed their entry into the conflict the model is able to accurately predict the group's tendency towards initiating conflict. On the other hand, actors like the Fulani and Hausa militias have notably positive sender and receiver effects (and the Christian and Muslim Militias have notably negative effects). 

% \begin{figure}[ht]
% 	\centering
% 	\includegraphics[width=.8\textwidth]{abEst2}
% 	\caption{Sender and receiver random effect estimates.}
% 	\label{fig:abEst}
% \end{figure}
% \FloatBarrier

\clearpage
\subsection{Network Goodness of Fit Assessment}

Figure~\ref{netGOF} presents an assessment of how well our model captures the network attributes of the conflict system in Nigeria on a variety of dimensions. For details on interpreting this diagnostic see \citet{minhas:etal:2016:arxiv}.

\begin{figure}[ht]
	\centering
	\includegraphics[width=1\textwidth]{netGOF}
	\caption{Network goodness of fit summary.}
	\label{netGOF}
	\end{figure}
\FloatBarrier
\clearpage

% \subsection{Rebel Group Information}

% stuff

\clearpage
\subsection{Temporal Forecasting Performance}

Scholars of violence are often concerned with models abilities to generated forecasts of future events. To do this, we compare the AME and GLM models' abilities to predict battles in the final year of our data with those in the rest of the data.

\begin{figure}[ht]
	\centering
	\begin{tabular}{cc}
	\includegraphics[width=.45\textwidth]{roc_outSampleForecast_1} & 
	\includegraphics[width=.45\textwidth]{rocPr_outSampleForecast_1}
	\end{tabular}
	\caption{Assessments of out-of-sample predictive performance when forecasting out the last period using ROC curves, separation plots, and PR curves. AUC statistics are provided as well for both curves. In each curve, the AME model is in blue, the GLM with lagged DV and covariates is in green, the GLM with only covariates is in red.}
	\label{fig:roc_ameForecast}
\end{figure}

\clearpage
\subsection{Additive and Multiplicative Effects Gibbs Sampler}

To estimate, the effects of our exogenous variables and latent attributes we utilize a Bayesian probit model in which we sample from the posterior distribution of the full conditionals until convergence. Specifically, given observed data $\textbf{Y}$ and $\textbf{X}$ -- where $\textbf{X}$ is a design array that includes our sender, receiver, and dyadic covariates -- we estimate our network of binary ties using a probit framework where: $y_{ij,t} = 1(\theta_{ij,t}>0)$ and $\theta_{ij,t} = \bm\beta^{\top}\mathbf{X}_{ij,t} + a_{i} + b_{j} + \textbf{u}_{i}^{\top} \textbf{D} \textbf{v}_{j} + \epsilon_{ij}$. The derivation of the full conditionals is described in detail in \citet{hoff:2005} and \citet{hoff:2008}, thus here we only outline the Markov chain Monte Carlo (MCMC) algorithm for the AME model that we utilize in this paper.

\begin{itemize}
 \item Given initial values of $\{\bm\beta, \textbf{a}, \textbf{b}, \textbf{U}, \textbf{V}, \Sigma_{ab}, \rho, \text{ and } \sigma_{\epsilon}^{2}\}$, the algorithm proceeds as follows:
 \begin{itemize}
 	\item sample $\bm\theta \; | \;  \bm\beta, \textbf{X}, \bm\theta, \textbf{a}, \textbf{b}, \textbf{U}, \textbf{V}, \Sigma_{ab}, \rho, \text{ and } \sigma_{\epsilon}^{2}$ (Normal)
 	\item sample $\bm\beta \; | \;  \textbf{X}, \bm\theta, \textbf{a}, \textbf{b}, \textbf{U}, \textbf{V}, \Sigma_{ab}, \rho, \text{ and } \sigma_{\epsilon}^{2}$ (Normal)
 	\item sample $\textbf{a}, \textbf{b} \; | \; \bm\beta, \textbf{X}, \bm\theta, \textbf{U}, \textbf{V}, \Sigma_{ab}, \rho, \text{ and } \sigma_{\epsilon}^{2}$ (Normal)
	\item sample $\Sigma_{ab} \; | \; \bm\beta, \textbf{X}, \bm\theta, \textbf{a}, \textbf{b}, \textbf{U}, \textbf{V}, \rho, \text{ and } \sigma_{\epsilon}^{2}$ (Inverse-Wishart)
 	\item update $\rho$ using a Metropolis-Hastings step with proposal $p^{*} | p  \sim$ truncated normal$_{[-1,1]}(\rho, \sigma_{\epsilon}^{2})$
 	\item sample $\sigma_{\epsilon}^{2} \; | \; \bm\beta, \textbf{X}, \bm\theta, \textbf{a}, \textbf{b}, \textbf{U}, \textbf{V}, \Sigma_{ab}, \text{ and } \rho$ (Inverse-Gamma)
 	\item For each $k \in K$:
 	\begin{itemize}
 		\item Sample $\textbf{U}_{[,k]} \; | \; \bm\beta, \textbf{X}, \bm\theta, \textbf{a}, \textbf{b}, \textbf{U}_{[,-k]}, \textbf{V}, \Sigma_{ab}, \rho, \text{ and } \sigma_{\epsilon}^{2}$ (Normal)
 		\item Sample $\textbf{V}_{[,k]} \; | \; \bm\beta, \textbf{X}, \bm\theta, \textbf{a}, \textbf{b}, \textbf{U}, \textbf{V}_{[,-k]}, \Sigma_{ab}, \rho, \text{ and } \sigma_{\epsilon}^{2}$ (Normal)
 		\item Sample $\textbf{D}_{[k,k]}  \; | \; \bm\beta, \textbf{X}, \bm\theta, \textbf{a}, \textbf{b}, \textbf{U}, \textbf{V}, \Sigma_{ab}, \rho, \text{ and } \sigma_{\epsilon}^{2}$ (Normal)\footnote{Subsequent to estimation, \textbf{D} matrix is absorbed into the calculation for $\textbf{V}$ as we iterate through $K$. }
 	\end{itemize}
 \end{itemize}
\end{itemize}

\newpage
\section{References}
\bibliography{/Users/cassydorff/ProjectsGit/whistle/master}
\bibliographystyle{jpr}
\end{document}
